% ============================================================
%  APPENDIX A : Supplementary Code
% ============================================================
\chapter{Supplementary Code: Simple Blockchain in Python}
\label{app:code}

\noindent\textbf{Declaration of Contributions:} This appendix was prepared by \textbf{Nitin Raj} and \textbf{Anmol Goklani}, who jointly wrote and tested the Python implementation.

\bigskip

To make the concepts discussed in this report concrete, we present a minimal implementation of a blockchain in Python. This implementation demonstrates the core ideas---hashing, block linking, and proof of work---in approximately 80 lines of code.

\begin{lstlisting}[language=Python, caption={A minimal blockchain implementation in Python.}, label={lst:blockchain}]
import hashlib
import json
import time

class Block:
    """Represents a single block in the blockchain."""

    def __init__(self, index, transactions, previous_hash, nonce=0):
        self.index = index
        self.timestamp = time.time()
        self.transactions = transactions
        self.previous_hash = previous_hash
        self.nonce = nonce
        self.hash = self.compute_hash()

    def compute_hash(self):
        block_string = json.dumps({
            "index": self.index,
            "timestamp": self.timestamp,
            "transactions": self.transactions,
            "previous_hash": self.previous_hash,
            "nonce": self.nonce
        }, sort_keys=True)
        return hashlib.sha256(block_string.encode()).hexdigest()


class Blockchain:
    """A simple blockchain with Proof of Work."""

    DIFFICULTY = 4  # number of leading zeros required

    def __init__(self):
        self.chain = []
        self.pending_transactions = []
        self.create_genesis_block()

    def create_genesis_block(self):
        genesis = Block(0, [], "0")
        self.chain.append(genesis)

    @property
    def last_block(self):
        return self.chain[-1]

    def proof_of_work(self, block):
        """Find a nonce such that the hash starts with
        DIFFICULTY number of leading zeros."""
        target = "0" * self.DIFFICULTY
        while not block.hash.startswith(target):
            block.nonce += 1
            block.hash = block.compute_hash()
        return block.hash

    def add_transaction(self, sender, receiver, amount):
        self.pending_transactions.append({
            "sender": sender,
            "receiver": receiver,
            "amount": amount
        })

    def mine_block(self):
        if not self.pending_transactions:
            print("No transactions to mine.")
            return False
        new_block = Block(
            index=self.last_block.index + 1,
            transactions=self.pending_transactions,
            previous_hash=self.last_block.hash
        )
        self.proof_of_work(new_block)
        self.chain.append(new_block)
        self.pending_transactions = []
        return new_block

    def is_chain_valid(self):
        for i in range(1, len(self.chain)):
            current = self.chain[i]
            previous = self.chain[i - 1]
            if current.previous_hash != previous.hash:
                return False
        return True


# --- Example Usage ---
if __name__ == "__main__":
    bc = Blockchain()
    bc.add_transaction("Alice", "Bob", 5.0)
    bc.add_transaction("Bob", "Charlie", 2.5)
    block = bc.mine_block()
    print(f"Block {block.index} mined: {block.hash}")
    print(f"Chain valid: {bc.is_chain_valid()}")
\end{lstlisting}

\paragraph{How to run.} Save the code as \texttt{blockchain.py} and execute:
\begin{lstlisting}[language=bash, numbers=none, frame=none]
$ python3 blockchain.py
Block 1 mined: 0000a7c3...  (hash with 4 leading zeros)
Chain valid: True
\end{lstlisting}

The \texttt{DIFFICULTY} parameter controls how many leading zeros the hash must have. Increasing it makes mining exponentially harder, illustrating the security--performance trade-off of Proof of Work.
